\subsection*{Общий случай принадлежности точки многоугольнику}Рассмотрим случайный луч, пущенный из точки. Если он пересек нечетное число сторон, то точка внутри.\newline
\textit{Что делать, если попали в вершину?} Пускать новый случайный луч, пока не повезет. Шанс попасть в вершину достаточно мал.

\subsection*{Множество запросов для выпуклого}
\Alg 
\begin{enumerate}
    \item Выберем самую нижнюю (если таких несколько, то самую левую среди них) вершину, теперь будем считать, что с нее начинается обход. Рассмотрим полярные углы, то есть углы между $\Vec{P_1P_i}$ и между положительным направлением оси $OX$.
    \item Заметим, что эти полярные углы отсортированы (по косинусу), значит по ним можно делать бинпоиск!
    \item Тогда, для того, чтобы проверить, лежит ли точка в многоугольнике, нужно посчитать угол между $\Vec{P_1A}$ и $OX$ и с помощью бинпоиска по полярным углам вершин получить, между какими векторами заключена точка(если она вообще оказалась заключена). Пусть это будут вектора $\Vec{P_1P_i}$ и $\Vec{P_1P_{i+1}}$. Тогда, если $S(\triangle P_1AP_i)+ S(\triangle P_1AP_i+1) \Le S(\triangle P_1P_iP_{i+1})$, то точка лежит внутри выпуклого многоугольника.

    Иначе, если $|[\Vec{P_1P_{i}}, \Vec{P_1A}]| + |[\Vec{P_1A}, \Vec{P_1P_{i+1}}]| \Le |[\Vec{P_1P_{i}}, \Vec{P_1P_{i+1}}]|$ и полярный угол $\Vec{P_1A}$ можно найти бинпоиском.
\end{enumerate}

\textbf{Асимптотика.} $\mathcal{O}(n + q\log(n))$, где $q$ - кол-во запросов.
\pagebreak